
\chapter{Base de données}
\label{chap:basededonnees}

\section{Modèle conceptuel}

La base de données d'AquaVolt est conçue pour centraliser les informations nécessaires à la gestion des clients, des contrats, des consommations d'eau et d'électricité, des factures et des opérations associées.

Elle repose sur plusieurs entités liées entre elles par des relations permettant de garantir la cohérence des informations.

Les principales entités du système sont :

\begin{itemize}[leftmargin=*]
    \item \textbf{Fournisseur :} représente l'organisme responsable de la gestion des clients et des tarifs.
    \item \textbf{Agent :} représente l'utilisateur chargé de l'importation des consommations annuelles.
    \item \textbf{Client :} représente l'abonné utilisant les services d'eau et d'électricité.
    \item \textbf{Quartier :} représente une zone géographique caractérisée par un type d'activité.
    \item \textbf{Contrat :} représente l'engagement entre un client et le service associé à un quartier.
    \item \textbf{Tarif :} définit les prix unitaires applicables selon le type de consommation et les tranches.
    \item \textbf{Consommation :} enregistre les index mensuels d'eau et d'électricité.
    \item \textbf{Photo\_capteur :} conserve les chemins des photographies associées aux relevés.
    \item \textbf{Anomalie :} enregistre les incohérences détectées dans les relevés.
    \item \textbf{Facture :} conserve les informations relatives à la facturation et au paiement.
    \item \textbf{Réclamation :} enregistre les demandes et signalements déposés par les clients.
    \item \textbf{Notification :} conserve les messages destinés aux différents utilisateurs.
    \item \textbf{Conso\_annuelle :} conserve les consommations annuelles importées par l'Agent.
\end{itemize}

Les relations entre ces entités permettent de représenter les différents processus métier de l'application.

\section{Modèle logique}

Le modèle logique est implémenté dans une base de données relationnelle nommée \texttt{miniproject\_version\_alpha}. Il utilise le moteur de stockage InnoDB et l'encodage UTF-8 avec le jeu de caractères \texttt{utf8mb4}.

Chaque table possède une clé primaire permettant d'identifier ses enregistrements. Les relations entre les tables sont établies à l'aide de clés étrangères.

Les principales associations sont les suivantes :

\begin{itemize}[leftmargin=*]
    \item Un fournisseur possède plusieurs clients et plusieurs tarifs.
    \item Un client peut posséder plusieurs contrats et plusieurs consommations.
    \item Un quartier peut être associé à plusieurs contrats.
    \item Un contrat possède plusieurs relevés mensuels.
    \item Une consommation peut être associée à des photographies et à des anomalies.
    \item Une consommation peut donner lieu à une facture unique.
    \item Un client peut déposer plusieurs réclamations.
    \item Un agent peut importer plusieurs relevés annuels.
    \item Un utilisateur peut recevoir plusieurs notifications.
\end{itemize}

\section{Description des tables}

La base de données comprend treize tables principales. Les sections suivantes présentent leur rôle et leurs principaux attributs.

\subsection{fournisseur}

La table \texttt{fournisseur} contient les informations relatives au fournisseur responsable de la gestion de la plateforme.

\begin{itemize}[leftmargin=*]
    \item \texttt{id\_fournisseur} : identifiant unique du fournisseur.
    \item \texttt{nom} : nom du fournisseur.
    \item \texttt{email} : adresse électronique, unique.
    \item \texttt{mot\_de\_passe} : mot de passe enregistré.
    \item \texttt{seuil\_eau} : seuil maximal défini pour la consommation d'eau.
    \item \texttt{seuil\_elec} : seuil maximal défini pour la consommation d'électricité.
    \item \texttt{date\_creation} : date de création de l'enregistrement.
\end{itemize}

\subsection{client}

La table \texttt{client} regroupe les informations personnelles et les informations d'accès des clients.

\begin{itemize}[leftmargin=*]
    \item \texttt{id\_client} : identifiant unique du client.
    \item \texttt{id\_fournisseur} : référence au fournisseur associé.
    \item \texttt{cin} : numéro de la carte d'identité nationale.
    \item \texttt{nom}, \texttt{prenom} : identité du client.
    \item \texttt{adresse} : adresse du client.
    \item \texttt{email} : adresse électronique unique.
    \item \texttt{mot\_de\_passe} : mot de passe enregistré.
    \item \texttt{telephone} : numéro de téléphone facultatif.
    \item \texttt{date\_creation} : date de création du compte.
\end{itemize}

\subsection{quartier}

La table \texttt{quartier} représente les quartiers pris en charge par l'application.

\begin{itemize}[leftmargin=*]
    \item \texttt{id\_quartier} : identifiant unique du quartier.
    \item \texttt{nom} : nom du quartier.
    \item \texttt{type\_activite} : type d'activité (\texttt{RESIDENTIEL}, \texttt{COMMERCIAL} ou \texttt{INDUSTRIEL}).
    \item \texttt{date\_creation} : date de création.
\end{itemize}

L'association entre le nom du quartier et son type d'activité doit être unique.

\subsection{contrat}

La table \texttt{contrat} permet de représenter les contrats associés aux clients et aux quartiers.

\begin{itemize}[leftmargin=*]
    \item \texttt{id\_contrat} : identifiant unique du contrat.
    \item \texttt{id\_client} : référence au client concerné.
    \item \texttt{id\_quartier} : référence au quartier associé.
    \item \texttt{numero\_contrat} : numéro unique du contrat.
    \item \texttt{date\_debut} : date de début du contrat.
    \item \texttt{date\_fin} : date de fin éventuelle.
    \item \texttt{statut} : état du contrat (\texttt{ACTIF} ou \texttt{RESILIE}).
    \item \texttt{date\_creation} : date de création de l'enregistrement.
\end{itemize}

\subsection{tarif}

La table \texttt{tarif} contient les tarifs appliqués par le fournisseur pour les consommations d'eau et d'électricité.

\begin{itemize}[leftmargin=*]
    \item \texttt{id\_tarif} : identifiant unique du tarif.
    \item \texttt{id\_fournisseur} : référence au fournisseur.
    \item \texttt{type} : type de consommation (\texttt{EAU} ou \texttt{ELEC}).
    \item \texttt{borne\_min} : limite inférieure de la tranche.
    \item \texttt{borne\_max} : limite supérieure de la tranche, éventuellement nulle pour la dernière tranche.
    \item \texttt{prix\_unitaire} : prix appliqué à l'unité consommée.
\end{itemize}

Les tranches sont représentées par des intervalles de la forme
\texttt{[borne\_min, borne\_max[}. Une borne maximale nulle indique une dernière tranche sans limite supérieure.

\subsection{consommation}

La table \texttt{consommation} enregistre les relevés mensuels des compteurs d'eau et d'électricité.

\begin{itemize}[leftmargin=*]
    \item \texttt{id\_conso} : identifiant unique du relevé.
    \item \texttt{id\_client} : référence au client.
    \item \texttt{id\_contrat} : référence au contrat concerné.
    \item \texttt{periode\_mois} : mois du relevé.
    \item \texttt{periode\_annee} : année du relevé.
    \item \texttt{index\_eau} : index du compteur d'eau.
    \item \texttt{index\_elec} : index du compteur d'électricité.
    \item \texttt{date\_saisie} : date d'enregistrement du relevé.
    \item \texttt{statut} : état du relevé.
\end{itemize}

Un contrat ne peut posséder qu'un seul relevé pour une même période mensuelle.

\subsection{photo\_capteur}

La table \texttt{photo\_capteur} conserve les références des photographies associées aux relevés des compteurs.

\begin{itemize}[leftmargin=*]
    \item \texttt{id\_photo} : identifiant unique de la photographie.
    \item \texttt{id\_conso} : référence à la consommation concernée.
    \item \texttt{type} : type de compteur (\texttt{EAU} ou \texttt{ELEC}).
    \item \texttt{chemin} : chemin du fichier photographique.
    \item \texttt{date\_televersement} : date du téléversement.
\end{itemize}

La contrainte d'unicité sur \texttt{id\_conso} et \texttt{type} limite à une seule photographie par type de compteur pour chaque consommation.

\subsection{anomalie}

La table \texttt{anomalie} permet d'enregistrer les incohérences détectées lors de la vérification des relevés.

\begin{itemize}[leftmargin=*]
    \item \texttt{id\_anomalie} : identifiant unique de l'anomalie.
    \item \texttt{id\_conso} : référence au relevé concerné.
    \item \texttt{id\_fournisseur} : référence au fournisseur chargé du traitement.
    \item \texttt{type\_anomalie} : type d'anomalie (\texttt{INDEX\_INFERIEUR} ou \texttt{SEUIL\_DEPASSE}).
    \item \texttt{champ} : compteur concerné (\texttt{EAU} ou \texttt{ELEC}).
    \item \texttt{valeur\_saisie} : valeur initialement saisie.
    \item \texttt{valeur\_reference} : valeur utilisée comme référence.
    \item \texttt{statut} : état de traitement de l'anomalie.
    \item \texttt{valeur\_corrigee} : valeur corrigée éventuelle.
    \item \texttt{date\_detection} : date de détection.
    \item \texttt{date\_resolution} : date de résolution éventuelle.
\end{itemize}

\subsection{facture}

La table \texttt{facture} contient les informations relatives aux factures générées à partir des consommations.

\begin{itemize}[leftmargin=*]
    \item \texttt{id\_facture} : identifiant unique de la facture.
    \item \texttt{numero} : numéro unique de la facture.
    \item \texttt{id\_client} : référence au client.
    \item \texttt{id\_conso} : référence à la consommation facturée.
    \item \texttt{date\_facture} : date d'émission.
    \item \texttt{periode\_mois} et \texttt{periode\_annee} : période facturée.
    \item \texttt{conso\_eau} et \texttt{conso\_elec} : consommations facturées.
    \item \texttt{montant\_ht} : montant hors taxes.
    \item \texttt{taux\_tva} : taux de TVA, fixé par défaut à 18\,\%.
    \item \texttt{montant\_tva} : montant de la TVA.
    \item \texttt{montant\_ttc} : montant toutes taxes comprises.
    \item \texttt{etat} : état de la facture (\texttt{IMPAYEE} ou \texttt{PAYEE}).
    \item \texttt{date\_paiement} : date du paiement éventuel.
    \item \texttt{mode\_paiement} : mode de paiement.
    \item \texttt{reference\_paiement} : référence du paiement.
\end{itemize}

La table permet ainsi de conserver les informations de facturation et de paiement sans utiliser une table de paiement distincte.

\subsection{notification}

La table \texttt{notification} conserve les notifications internes destinées aux différents utilisateurs de l'application.

\begin{itemize}[leftmargin=*]
    \item \texttt{id\_notification} : identifiant unique de la notification.
    \item \texttt{destinataire\_role} : rôle du destinataire (\texttt{client}, \texttt{fournisseur} ou \texttt{agent}).
    \item \texttt{id\_destinataire} : identifiant du destinataire.
    \item \texttt{type} : catégorie de notification.
    \item \texttt{titre} : titre du message.
    \item \texttt{message} : contenu de la notification.
    \item \texttt{lien} : lien facultatif vers une fonctionnalité.
    \item \texttt{lu} : indicateur de lecture.
    \item \texttt{date\_creation} : date de création.
    \item \texttt{date\_lecture} : date de lecture éventuelle.
\end{itemize}

Le destinataire est identifié par la combinaison de son rôle et de son identifiant. Cette conception permet de gérer plusieurs catégories d'utilisateurs dans une même table.

\subsection{agent}

La table \texttt{agent} contient les informations relatives aux agents chargés notamment de l'importation annuelle.

\begin{itemize}[leftmargin=*]
    \item \texttt{id\_agent} : identifiant unique de l'agent.
    \item \texttt{nom} et \texttt{prenom} : identité de l'agent.
    \item \texttt{email} : adresse électronique unique.
    \item \texttt{mot\_de\_passe} : mot de passe enregistré.
    \item \texttt{telephone} : numéro de téléphone facultatif.
    \item \texttt{date\_creation} : date de création du compte.
\end{itemize}

\subsection{conso\_annuelle}

La table \texttt{conso\_annuelle} conserve les relevés annuels issus de la tournée de l'Agent.

\begin{itemize}[leftmargin=*]
    \item \texttt{id\_conso\_annuelle} : identifiant unique du relevé annuel.
    \item \texttt{id\_client} : référence au client concerné.
    \item \texttt{id\_agent} : référence à l'agent ayant effectué l'importation.
    \item \texttt{conso\_eau} : consommation annuelle d'eau.
    \item \texttt{conso\_elec} : consommation annuelle d'électricité.
    \item \texttt{annee} : année concernée.
    \item \texttt{date\_saisie} : date d'enregistrement.
    \item \texttt{statut\_verification} : état de vérification.
    \item \texttt{date\_verification} : date de vérification éventuelle.
\end{itemize}

La combinaison du client et de l'année est unique, ce qui empêche l'enregistrement de plusieurs relevés annuels pour un même client durant la même année.

\subsection{reclamation}

La table \texttt{reclamation} permet de gérer les réclamations déposées par les clients.

\begin{itemize}[leftmargin=*]
    \item \texttt{id\_reclamation} : identifiant unique de la réclamation.
    \item \texttt{numero} : numéro unique de la réclamation.
    \item \texttt{id\_client} : référence au client concerné.
    \item \texttt{id\_facture} : référence facultative à une facture.
    \item \texttt{type} : type de réclamation.
    \item \texttt{description} : description du problème signalé.
    \item \texttt{piece\_jointe} : chemin d'une pièce jointe éventuelle.
    \item \texttt{statut} : état de traitement.
    \item \texttt{reponse} : réponse éventuelle du fournisseur.
    \item \texttt{date\_creation} : date de dépôt.
    \item \texttt{date\_resolution} : date de résolution éventuelle.
\end{itemize}

\section{Relations et cardinalités}

Les relations entre les tables sont définies à l'aide de clés étrangères. Les principales cardinalités du modèle sont présentées dans le tableau suivant.

\begin{table}[H]
    \centering
    \caption{Principales relations entre les entités}
    \begin{tabularx}{\textwidth}{|l|l|X|}
        \hline
        \textbf{Entités liées} & \textbf{Cardinalité} & \textbf{Signification} \\
        \hline
        Fournisseur -- Client & 1,N & Un fournisseur peut gérer plusieurs clients. \\
        \hline
        Fournisseur -- Tarif & 1,N & Un fournisseur peut définir plusieurs tarifs. \\
        \hline
        Client -- Contrat & 1,N & Un client peut posséder plusieurs contrats. \\
        \hline
        Quartier -- Contrat & 1,N & Un quartier peut être associé à plusieurs contrats. \\
        \hline
        Client -- Consommation & 1,N & Un client peut posséder plusieurs relevés. \\
        \hline
        Contrat -- Consommation & 1,N & Un contrat peut posséder plusieurs relevés mensuels. \\
        \hline
        Consommation -- Photo\_capteur & 1,N & Un relevé peut être associé à plusieurs photographies. \\
        \hline
        Consommation -- Anomalie & 1,N & Un relevé peut présenter plusieurs anomalies. \\
        \hline
        Consommation -- Facture & 1,1 & Une consommation peut être associée à une seule facture au maximum. \\
        \hline
        Client -- Facture & 1,N & Un client peut recevoir plusieurs factures. \\
        \hline
        Client -- Réclamation & 1,N & Un client peut déposer plusieurs réclamations. \\
        \hline
        Agent -- Conso\_annuelle & 1,N & Un agent peut être associé à plusieurs relevés annuels. \\
        \hline
        Client -- Conso\_annuelle & 1,N & Un client peut avoir plusieurs relevés annuels, un par année. \\
        \hline
    \end{tabularx}
\end{table}

La relation entre une facture et une réclamation est facultative du côté de la réclamation : une réclamation peut concerner une facture précise ou être indépendante d'une facture.

La table \texttt{notification} constitue un cas particulier. Elle identifie son destinataire à l'aide de deux attributs, \texttt{destinataire\_role} et \texttt{id\_destinataire}, sans clé étrangère directe vers les tables \texttt{client}, \texttt{fournisseur} ou \texttt{agent}.

\section{Contraintes d'intégrité}

Plusieurs contraintes permettent de préserver la cohérence des données.

\subsection*{Clés primaires}

Chaque table possède une clé primaire auto-incrémentée de type \texttt{INT UNSIGNED}, permettant d'identifier chaque enregistrement de manière unique.

\subsection*{Clés étrangères}

Les clés étrangères assurent la cohérence des relations entre les tables. Elles utilisent principalement les règles suivantes :

\begin{itemize}[leftmargin=*]
    \item \texttt{ON UPDATE CASCADE} : les modifications d'identifiants référencés sont propagées.
    \item \texttt{ON DELETE RESTRICT} : la suppression d'un enregistrement référencé est empêchée.
    \item \texttt{ON DELETE CASCADE} : la suppression d'un client entraîne la suppression de ses relevés annuels ; la suppression d'une consommation entraîne celle de ses photographies et de ses anomalies.
    \item \texttt{ON DELETE SET NULL} : la référence à une facture est mise à \texttt{NULL} lorsqu'une facture associée est supprimée.
\end{itemize}

\subsection*{Contraintes CHECK}

Des contraintes de vérification sont définies pour empêcher l'enregistrement de valeurs incohérentes :

\begin{itemize}[leftmargin=*]
    \item les seuils de consommation doivent être positifs ou nuls ;
    \item les bornes tarifaires doivent être cohérentes et les prix unitaires positifs ou nuls ;
    \item les dates de fin des contrats ne peuvent pas être antérieures aux dates de début ;
    \item les périodes mensuelles doivent être comprises entre 1 et 12 ;
    \item les index et les consommations doivent être positifs ou nuls ;
    \item les années des relevés annuels doivent être comprises entre 2000 et 2100 ;
    \item les montants enregistrés dans les factures doivent être positifs ou nuls.
\end{itemize}

\section{Index et contraintes d'unicité}

Afin de garantir l'unicité de certaines informations et de faciliter les recherches, plusieurs index sont définis dans la base de données.

Les principales contraintes d'unicité sont les suivantes :

\begin{itemize}[leftmargin=*]
    \item adresse électronique du fournisseur ;
    \item adresse électronique de l'agent ;
    \item CIN et adresse électronique du client ;
    \item association du nom et du type d'activité du quartier ;
    \item numéro du contrat ;
    \item combinaison du fournisseur, du type de tarif et de la borne minimale ;
    \item combinaison du contrat, du mois et de l'année pour une consommation ;
    \item combinaison du client et de l'année pour une consommation annuelle ;
    \item combinaison de la consommation et du type de photographie ;
    \item numéro de facture ;
    \item référence de consommation dans la table des factures ;
    \item numéro de réclamation.
\end{itemize}

Des index supplémentaires sont également définis sur les clés étrangères et certains attributs utilisés pour les recherches, notamment les destinataires des notifications et leurs dates de création.

\section{États et transitions}

Plusieurs tables utilisent des attributs de type \texttt{ENUM} afin de représenter les différents états des opérations métier.

\subsection*{État d'un contrat}

Un contrat peut prendre deux états :

\begin{itemize}[leftmargin=*]
    \item \texttt{ACTIF} : le contrat est en cours.
    \item \texttt{RESILIE} : le contrat a été résilié.
\end{itemize}

\subsection*{État d'une consommation}

Une consommation peut prendre les états suivants :

\begin{itemize}[leftmargin=*]
    \item \texttt{EN\_ATTENTE} : le relevé vient d'être enregistré.
    \item \texttt{ANOMALIE} : une incohérence a été détectée.
    \item \texttt{VALIDEE} : le relevé a été validé.
    \item \texttt{FACTUREE} : une facture a été générée pour le relevé.
\end{itemize}

\subsection*{État d'une anomalie}

Le traitement d'une anomalie utilise trois états :

\begin{itemize}[leftmargin=*]
    \item \texttt{EN\_ATTENTE} : anomalie en attente de traitement.
    \item \texttt{CONFIRMEE} : anomalie confirmée.
    \item \texttt{CORRIGEE} : anomalie corrigée.
\end{itemize}

\subsection*{État d'une facture}

Une facture peut être :

\begin{itemize}[leftmargin=*]
    \item \texttt{IMPAYEE} : le paiement n'a pas encore été enregistré.
    \item \texttt{PAYEE} : le paiement a été enregistré.
\end{itemize}

\subsection*{État d'une réclamation}

Une réclamation peut prendre les états suivants :

\begin{itemize}[leftmargin=*]
    \item \texttt{OUVERTE} : la réclamation vient d'être déposée.
    \item \texttt{EN\_COURS} : son traitement est en cours.
    \item \texttt{RESOLUE} : le problème a été résolu.
    \item \texttt{REJETEE} : la réclamation a été rejetée.
\end{itemize}

\subsection*{État de vérification annuelle}

Les relevés annuels utilisent trois états :

\begin{itemize}[leftmargin=*]
    \item \texttt{A\_VERIFIER} : relevé en attente de vérification.
    \item \texttt{VALIDEE} : relevé validé.
    \item \texttt{REJETEE} : relevé rejeté.
\end{itemize}

\subsection*{État des notifications}

L'attribut \texttt{lu} indique si une notification a été consultée. La valeur 0 signifie qu'elle n'a pas encore été lue, tandis que la valeur 1 indique qu'elle a été lue.

\section{Jeu de données de démonstration}

Le projet prévoit un fichier \texttt{seed.sql} destiné à préparer les données initiales de la base de données pour les besoins de démonstration et de test.

Ce fichier est distinct du script \texttt{schema.sql}, qui définit la structure des tables, leurs attributs, leurs clés et leurs contraintes.

\section{Migrations}

Le projet comporte un répertoire \texttt{sql/migrations/} destiné à regrouper les scripts de migration de la base de données.

Le script \texttt{schema.sql} constitue le script de création initiale de la structure. Il crée la base \texttt{miniproject\_version\_alpha}, définit l'encodage utilisé et crée les treize tables dans un ordre compatible avec leurs dépendances.

Les scripts de migration permettent de conserver séparément les éventuelles évolutions de la structure de la base de données au cours du développement.
